\section{Discussion, Limitations, and Conclusion}\label{sec:conclusion}

\subsection{Computational efficiency}

Across the verification grid the hand-written solver required between \textbf{222 and 2113}
iterations. Per-iteration cost is dominated by a single matrix--vector product $\Sig w$ at
$O(N^2)$, plus the bisection of \cref{prop:jointprox} at $O(N \log(1/\epsilon))$; no matrix
factorization occurs anywhere inside the loop, and memory traffic is one $N \times N$ read per
iteration.

The contrast with the interior-point method is the expected one in kind: CLARABEL converges in
a few dozen iterations rather than hundreds, but each iteration forms and factorizes a KKT
system at roughly $O(N^3)$. It would be convenient to conclude that the first-order method
therefore wins, and the data does not support that. At $N = 98$ a dense $98 \times 98$
factorization is trivial for a modern machine, the interior-point solver is entirely practical,
and neither method is anywhere near a wall-clock bottleneck. The first-order method's advantage
here is that it is transparent and was built from the problem structure rather than invoked;
its asymptotic advantage would become a real one only at $N$ in the thousands, where
factorization cost dominates and the $O(N^2)$ per-iteration profile starts to matter.

\subsection{Solution quality}

The solver produces exact zeros rather than thresholded small values, satisfies the budget
constraint to $10^{-12}$, matches an independently implemented interior-point solver to a
maximum relative objective gap of $7.75 \times 10^{-6}$, and reproduces its active set exactly
on every parameter set tested (\cref{tab:cvxpy}). For a convex problem with a unique optimal
value, that combination is strong evidence of correctness.

\subsection{Limitations}

\begin{itemize}
  \item \textbf{Survivorship bias.} The current VN100 membership is applied backward over four
        years, so delisted and removed tickers are absent from the universe.
  \item \textbf{$\muhat$ remains a poor estimator.} The sample mean has high variance over a
        short window. The robust term blunts the consequences of this; it does not remove them,
        and no term in \cref{eq:main} can.
  \item \textbf{The $\kappa$ finding is empirical.} The increase in concentration reported in
        \cref{sec:insample} is an observation on one covariance structure with a plausible
        mechanism, not a general result.
  \item \textbf{The out-of-sample window is short.} Roughly 24 months spanning a single market
        regime. The conclusions of \cref{sec:backtest} cannot be extrapolated to other cycles.
  \item \textbf{Validation overfitting.} Selecting $(\kappa,\gamma)$ on a six-month window is a
        mild form of fitting to the validation set, inherent to short-window walk-forward
        rather than a defect of this implementation --- and, as \cref{sec:backtest} shows, it
        is the direct cause of the turnover that drives the cost disadvantage.
  \item \textbf{A single historical path.} No block bootstrap or resampling was performed, so
        no confidence interval accompanies the reported Sharpe ratio or drawdown.
  \item \textbf{Costs are modelled crudely.} A flat $0.20\%$ of turnover, with no market-impact
        or slippage term and no dependence on order size or liquidity.
\end{itemize}

\subsection{Conclusion}

The portfolio selection problem was formulated as a convex, second-order-cone representable
program in which the robust term is derived as an exact worst-case objective over an
ellipsoidal uncertainty set rather than assumed, and in which the $\ell_1$ penalty serves
simultaneously as a convex surrogate for the number of holdings and as a cap on gross leverage.
A proximal-subgradient method was implemented in pure NumPy, resting on an exact joint
proximal operator for the $\ell_1$ penalty and the budget constraint taken together
(\cref{prop:jointprox}). That solver reproduces an independent interior-point solver's optimal
value to $7.75 \times 10^{-6}$ relative gap and its active set exactly, on every parameter set
tested. Out of sample, the resulting long-only strategy beat an equal-weight benchmark on a
risk-adjusted basis over 494 trading days --- by a real but modest margin, at materially higher
volatility, deeper drawdown, and roughly seven times the transaction cost.

The methodological lesson is the one from \cref{sec:proxbug}. The original solver composed two
individually correct operations --- soft thresholding and projection onto the budget hyperplane
--- into an incorrect one, and returned portfolios that were feasible, plausible, and wrong by
several percent in objective and by a factor of five in the number of positions held. Nothing
in the output revealed it. It was found only by solving the same problem a second time with an
unrelated algorithm and comparing. For convex problems, where a unique optimum makes such a
comparison decisive, building that independent check is worth more than any amount of
inspection of the primary implementation.

\label{sec:endofbody}
